\documentclass[11pt]{article}
\usepackage{amsmath, amssymb, amsthm}
\usepackage[margin=1in]{geometry}
\usepackage{hyperref}

\newtheorem{theorem}{Theorem}
\newtheorem{remark}{Remark}

\title{Anonymized Evaluation with Statistical Reliability:\\Krippendorff's Alpha, Calibration, and Score Aggregation}
\author{Abhi Wadhwa}
\date{}

\begin{document}
\maketitle

\begin{abstract}
Notes on the mathematical foundations behind a double-blind review platform. The main technical content is Krippendorff's alpha for measuring inter-rater reliability, which I derive in some detail because the standard references make it more confusing than it needs to be. I also cover score aggregation methods (weighted average, trimmed mean, median), outlier detection, reviewer calibration, and the balanced assignment problem. This is essentially the math you need if you want to build a fair, statistically grounded peer review system.
\end{abstract}

\section{Why Krippendorff's Alpha?}

The central question in any peer review system is: are the reviewers actually agreeing with each other, or are they just assigning random numbers? There are a bunch of inter-rater reliability coefficients out there --- Cohen's kappa, Fleiss' kappa, intraclass correlation --- but I went with Krippendorff's alpha because it handles \textbf{any number of raters, missing data, and multiple measurement levels} (nominal, ordinal, interval, ratio) all in one framework. It's the Swiss army knife of agreement measures.

The core formula is deceptively simple:
\[
    \alpha = 1 - \frac{D_o}{D_e}
\]
where $D_o$ is the observed disagreement among raters and $D_e$ is the expected disagreement if ratings were assigned by chance. If raters agree perfectly, $D_o = 0$ and $\alpha = 1$. If they agree only as much as chance would predict, $D_o = D_e$ and $\alpha = 0$. If they somehow disagree \textit{more} than chance (systematic opposition), $\alpha < 0$.

\section{Computing $D_o$ and $D_e$ for Interval Data}

I spent a while trying to figure out why different sources give slightly different formulas for $D_o$. It turns out the confusion comes from whether you're working with the ``reliability data matrix'' (units as rows, raters as columns, lots of missing cells) or the ``coincidence matrix'' approach. Krippendorff's own formulation uses the coincidence matrix, and that's what I implemented. Let me walk through it for interval data since that's what review scores are.

Suppose we have $N$ units (submissions) and for each unit $u$, there are $m_u$ raters who assigned scores. Let $x_{u,c}$ denote the score assigned by rater $c$ to unit $u$. The \textbf{observed disagreement} is:
\[
    D_o = \frac{1}{n} \sum_{u=1}^{N} \frac{1}{m_u - 1} \sum_{\substack{c, c' \\ c < c'}} (x_{u,c} - x_{u,c'})^2
\]
where $n = \sum_u \binom{m_u}{2}$ is the total number of pairable values across all units. Notice that the $1/(m_u - 1)$ factor is crucial --- it's what allows Krippendorff's alpha to handle different numbers of raters per unit. If every unit has exactly $m$ raters, this simplifies, but the general case handles missing data gracefully.

For the \textbf{expected disagreement}, we compute what $D_o$ would look like if all ratings were just drawn from the marginal distribution of all observed values. Let $\bar{n} = \sum_u m_u$ be the total number of observed ratings. Then:
\[
    D_e = \frac{1}{\bar{n}(\bar{n} - 1)} \sum_{\substack{(u,c), (u',c') \\ (u,c) \neq (u',c')}} (x_{u,c} - x_{u',c'})^2
\]
This is essentially the variance of all observed scores, scaled appropriately. In practice, you can compute it more efficiently as:
\[
    D_e = \frac{2}{\bar{n}(\bar{n}-1)} \left( \bar{n} \sum_{u,c} x_{u,c}^2 - \left(\sum_{u,c} x_{u,c}\right)^2 \right)
\]
which avoids the double sum.

\subsection{Interpretation}

The standard interpretation thresholds (from Krippendorff's own recommendations) are:

\begin{center}
\begin{tabular}{ll}
    $\alpha \geq 0.800$ & Reliable agreement \\
    $0.667 \leq \alpha < 0.800$ & Tentative --- draw cautious conclusions \\
    $\alpha < 0.667$ & Unreliable --- re-examine the review process \\
    $\alpha < 0$ & Systematic disagreement (raters are anti-correlated)
\end{tabular}
\end{center}

\textbf{The key insight} is that $\alpha$ corrects for chance agreement, unlike a raw percentage. Two raters who both assign 3/5 to everything will have high raw agreement but $\alpha$ will flag this as meaningless (no variance $\Rightarrow$ $D_e \approx 0$, making $\alpha$ undefined or unstable). This is exactly the behavior you want.

\section{Score Aggregation}

Once you've verified that reviewers are reasonably consistent (via $\alpha$), you need to aggregate multiple scores into a single ranking. I implemented three methods, each with its own strengths.

\subsection{Weighted Average}

The default. Each reviewer $i$ has a weight $w_i$ (which can come from the calibration step, see below), and the aggregated score for a submission on dimension $d$ is:
\[
    \bar{s}_d = \frac{\sum_{i} w_i \cdot s_{i,d}}{\sum_i w_i}
\]
This is just a weighted mean. Simple, interpretable, and it lets you down-weight reviewers who were inconsistent during calibration. If all weights are equal, it reduces to the arithmetic mean.

\subsection{Trimmed Mean}

Remove the top and bottom $k\%$ of scores before averaging. This is robust to outliers --- if one reviewer is consistently harsh or lenient, their extreme scores get dropped. You need at least 3 reviews per submission for this to make sense, because with 2 reviews, trimming one from each end leaves you with nothing.

Formally, given scores $s_{(1)} \leq s_{(2)} \leq \cdots \leq s_{(m)}$ for a submission, the $k\%$-trimmed mean is:
\[
    \bar{s}_{\text{trim}} = \frac{1}{m - 2\lfloor km/100 \rfloor} \sum_{i = \lfloor km/100 \rfloor + 1}^{m - \lfloor km/100 \rfloor} s_{(i)}
\]

\subsection{Median}

The most robust option. Completely ignores outliers (any single extreme score has no effect). The downside is reduced sensitivity --- small genuine differences between submissions get washed out.

\section{Outlier Detection}

Before aggregation, I flag individual scores that are more than $N_\sigma$ standard deviations from the mean for that submission and dimension:
\[
    |s_{i,d} - \bar{s}_d| > N_\sigma \cdot \sigma_d
\]
where $\bar{s}_d$ and $\sigma_d$ are computed from all scores on dimension $d$ for that submission. Flagged scores are excluded from the final aggregation. In practice, I use $N_\sigma = 2$ as the default threshold. After playing with this for a while, I found that $N_\sigma = 2$ catches genuinely anomalous scores without being so aggressive that it trims legitimate disagreement. Setting it to $1.5$ was too aggressive; $3$ was too lenient for the small sample sizes typical in peer review (3--5 reviewers per submission).

\section{Reviewer Calibration}

The calibration step is how you figure out which reviewers are trustworthy before the real reviews start. The idea is simple: have all reviewers score the same set of calibration applications, then measure how much each reviewer deviates from the consensus.

For reviewer $i$ on calibration application $a$ and dimension $d$, the deviation is:
\[
    e_{i,a,d} = |s_{i,a,d} - \bar{s}_{a,d}|
\]
where $\bar{s}_{a,d}$ is the mean score across all reviewers for that application and dimension. The reviewer's overall calibration score is their mean absolute deviation across all calibration items and dimensions:
\[
    \text{MAD}_i = \frac{1}{|A| \cdot |D|} \sum_{a \in A} \sum_{d \in D} e_{i,a,d}
\]
Reviewers with high MAD are flagged. You can either retrain them, remove them, or (what I do) use their MAD to compute a weight $w_i = 1 / (1 + \text{MAD}_i)$ for the weighted average aggregation. This way, inconsistent reviewers are automatically down-weighted rather than excluded entirely, which feels more fair.

\section{Balanced Assignment}

The assignment problem is: given $M$ submissions and $K$ reviewers, assign exactly $k$ reviewers to each submission such that (a) no reviewer has a conflict of interest with any assigned submission, and (b) reviewer workloads differ by at most 1.

The total number of review slots is $M \cdot k$. If divided evenly, each reviewer handles $\lfloor Mk / K \rfloor$ or $\lceil Mk / K \rceil$ reviews. I use a simple round-robin algorithm:

\begin{enumerate}
    \item Shuffle the reviewer list (for fairness).
    \item Maintain a pointer into the reviewer list.
    \item For each submission, assign the next $k$ non-conflicted reviewers from the pointer position, wrapping around as needed.
    \item If a conflict is detected, skip that reviewer and try the next.
\end{enumerate}

This is a greedy algorithm and it's not guaranteed to find a valid assignment if conflicts are very dense, but in practice with reasonable conflict rates (say, $<20\%$), it works well. For adversarial cases, you'd want to formulate this as a constraint satisfaction problem or an integer program, but I haven't needed that yet.

The workload balance guarantee comes from the round-robin structure: the pointer advances uniformly, so no reviewer can be more than 1 assignment ahead of any other. The ``at most 1'' difference is tight when $Mk$ is not divisible by $K$.

\section{Putting It All Together}

The full pipeline is: anonymize $\to$ assign $\to$ calibrate $\to$ review $\to$ detect outliers $\to$ aggregate $\to$ compute $\alpha$ $\to$ rank. Each step has clear mathematical justification, and the whole system is designed so that the final ranking is as unbiased and statistically reliable as possible. Krippendorff's alpha serves as the quality control metric --- if it's below 0.667, you know something is wrong with your review process and you should investigate before trusting the results.

\begin{thebibliography}{9}
\bibitem{kripp} Krippendorff, K., ``Content Analysis: An Introduction to Its Methodology,'' 4th ed., SAGE, 2018.
\bibitem{kripp2} Krippendorff, K., ``Computing Krippendorff's Alpha-Reliability,'' \url{https://repository.upenn.edu/asc_papers/43/}.
\bibitem{shrout} Shrout, P.E. and Fleiss, J.L., ``Intraclass Correlations: Uses in Assessing Rater Reliability,'' \textit{Psychological Bulletin}, 86(2), 1979.
\end{thebibliography}

\end{document}
